\hypertarget{reflections-on-point-image-phase-synchronization-with-commander-sol}{%
\section{Reflections on Point-Image Phase Synchronization with Commander
Sol}\label{reflections-on-point-image-phase-synchronization-with-commander-sol}}

\hypertarget{exact-costs-cut-space-packings-and-a-sharp-stabilization-threshold-in-fcoa}{%
\subsection{Exact Costs, Cut-Space Packings, and a Sharp Stabilization
Threshold in
FCOA}\label{exact-costs-cut-space-packings-and-a-sharp-stabilization-threshold-in-fcoa}}

\textbf{Alex Malachevsky · Commander Sol}\\
\textbf{Research manuscript · Version 1.0-rc1 · 31 August 2026}

\begin{center}\rule{0.5\linewidth}{0.5pt}\end{center}

\hypertarget{abstract}{%
\subsection{Abstract}\label{abstract}}

This paper works in the Fixed-Carrier Oriented Algebra (FCOA) framework
fixed by \textbf{FCOA Definition 1.0}, \emph{Fixed-Carrier Oriented
Algebra (FCOA): Definition, Typed Partial Operations, Carrier Erasure,
and the Canonical M0 Baseline}, https://doi.org/10.5281/zenodo.22164246.
We study an abstract full-support phase-synchronization problem that
arises after sparse anonymous operation components admit well-defined
local color phases. Given \texttt{r} permutations
\texttt{pi\_1,...,pi\_r\ in\ S\_q}, a primitive point-image constraint
has the form \texttt{pi\_i(a)=pi\_j(a)} for one source color \texttt{a}.
Let \texttt{L\_q(r)} be the minimum number of such constraints that
force every satisfying tuple to be diagonal.

We prove an exact graph-theoretic reformulation: a constraint system
synchronizes exactly when a canonical transversal quotient is uniquely
\texttt{q}-colorable up to global color relabeling. This yields the
necessary pair-union connectivity of every two source-color constraint
graphs. We then determine several infinite exact families,

\[
L_2(r)=r-1,
\qquad
L_3(r)=\left\lceil\frac{3(r-1)}2\right\rceil,
\qquad
L_q(2)=q-1,
\qquad
L_q(3)=2q-3,
\]

and close the entire four-phase column:

\[
L_q(4)=
\begin{cases}
3,&q=2,\\
2q-1,&3\le q\le5,\\
12,&q=6,\\
3q-7,&q\ge7.
\end{cases}
\]

The main structural result associates to every source color a normalized
binary cut space \texttt{W(P\_a)\ \textless{}=\ F\_2\^{}\{r-1\}} whose
dimension equals the saved constraint cost of that color.
Synchronization forces the nonzero parts of these cut spaces to be
pairwise disjoint, so

\[
\sum_a(2^{d_a}-1)\le2^{r-1}-1.
\]

A matching binary-cut construction then gives the exact large-alphabet
law

\[
\boxed{L_q(r)=(r-1)q-(2^{r-1}-1)}
\]

for every

\[
q\ge2^{r-1}-1,
\]

and we prove that this threshold is exact. Thus the unresolved part of
the problem is finite for every fixed \texttt{r}. The finite-geometry
ingredients themselves are classical; the contribution claimed here is
the FCOA point-image synchronization model, its unique-coloring and
cut-space reductions, the exact families above, and the sharp
stabilization theorem.

\textbf{Keywords:} FCOA; anonymous values; permutation synchronization;
unique colorability; set partitions; cut spaces; partial spreads;
rigidity cost.

\begin{center}\rule{0.5\linewidth}{0.5pt}\end{center}

\hypertarget{introduction}{%
\subsection{1. Introduction}\label{introduction}}

Anonymous-output layers in FCOA naturally separate two questions. The
first is whether an automorphism of a sparse relational reduct induces a
local permutation of visible output names. The second begins only after
such local phases exist: how much additional equality information is
required to force all local phases to agree with one global anonymous
relabeling?

The present paper isolates the second question as a finite extremal
problem.

Let

\[
O=\{0,1,\ldots,q-1\}
\]

be an anonymous alphabet and let

\[
\pi_1,\ldots,\pi_r\in S_O.
\]

For \texttt{1\textless{}=i\textless{}j\textless{}=r} and
\texttt{a\ in\ O}, a primitive \textbf{point-image constraint} is

\[
[i,j;a]:\qquad \pi_i(a)=\pi_j(a).
\tag{1}
\]

A set \texttt{S} of such constraints is \textbf{synchronizing} if every
satisfying tuple is diagonal:

\[
\pi_1=\cdots=\pi_r.
\tag{2}
\]

We define

\[
\boxed{
L_q(r)=\min\{|S|:S\text{ is synchronizing}\}.
}
\tag{3}
\]

The old spanning-tree construction immediately gives

\[
L_q(r)\le(q-1)(r-1),
\tag{4}
\]

because equality of two permutations on \texttt{q-1} source points
forces equality everywhere. For \texttt{q=2}, this is exact. For
\texttt{q\textgreater{}=3}, however, bijectivity allows constraints on
different component pairs to interact, and the tree bound is far from
the whole story.

The main results of this paper are:

\begin{enumerate}
\def\labelenumi{\arabic{enumi}.}
\tightlist
\item
  an exact reduction of synchronization to unique colorability of a
  special transversal quotient;
\item
  a universal pair-union connectivity condition and the lower bound \[
  L_q(r)\ge\left\lceil\frac{q(r-1)}2\right\rceil;
  \tag{5}
  \]
\item
  exact formulas on the complete \texttt{q=3} row and \texttt{r=2,3,4}
  columns;
\item
  a universal binary-cut construction using all nonzero vectors of
  \texttt{F\_2\^{}\{r-1\}};
\item
  a cut-space packing inequality that matches that construction and
  yields the exact stabilization theorem \[
  L_q(r)=(r-1)q-(2^{r-1}-1)
  \quad(q\ge2^{r-1}-1),
  \tag{6}
  \] with exact threshold.
\end{enumerate}

The finite pre-stabilization region remains open for
\texttt{r\textgreater{}=5}; this is deliberately separated from the
proved tail.

\begin{center}\rule{0.5\linewidth}{0.5pt}\end{center}

\hypertarget{fcoa-framework-and-scope}{%
\subsection{2. FCOA framework and
scope}\label{fcoa-framework-and-scope}}

\hypertarget{framework-statement}{%
\subsubsection{2.1 Framework statement}\label{framework-statement}}

The canonical framework is FCOA Definition 1.0 {[}1{]}. The underlying
FCOA motivation is a finite carrier \texttt{G} with a partial operation
whose terminal outputs form anonymous value fibers. Sparse relational
reducts can decompose the defined-cell geometry into comparison
components. In the \textbf{full-support phase-admissible sector}, each
relevant component carries a local permutation of the same anonymous
terminal alphabet \texttt{O}.

This paper passes to the derived synchronization layer. The phase-index
carrier is

\[
R=\{1,\ldots,r\},
\]

and the terminal sort is the anonymous set \texttt{O}. The primitive
data studied here are not new operation values or new real FCOA cells;
they are the abstract equality tests (1) between images of one source
color under two local phases. Relative to the preceding sparse
multicolor transport layer, the only added structure is a selected
family of these point-image equalities.

\hypertarget{erasure-and-recovery}{%
\subsubsection{2.2 Erasure and recovery}\label{erasure-and-recovery}}

The \textbf{erasure convention} is that names of terminal values remain
anonymous: simultaneous left composition

\[
\pi_i\mapsto\sigma\circ\pi_i
\qquad(\sigma\in S_q)
\tag{7}
\]

changes only the common output naming and preserves every constraint.
Consequently one may normalize one phase to the identity.

The \textbf{recovery target} is not recovery of named colors. It is
recovery of one global anonymous relabeling, equivalently
diagonalization (2).

\hypertarget{firewall}{%
\subsubsection{2.3 Firewall}\label{firewall}}

The quantity \texttt{L\_q(r)} is an abstract full-support
phase-synchronization cost. It is not the minimum number of real
operation cells needed to repair an FCOA carrier. No multicolor
real-cell invariant \texttt{alpha\_q} is defined here, and no inequality
comparing such an invariant with \texttt{L\_q(r)} is asserted. No
arithmetic on the FCOA carrier is imported.

\begin{center}\rule{0.5\linewidth}{0.5pt}\end{center}

\hypertarget{constraint-graphs-and-forest-reduction}{%
\subsection{3. Constraint graphs and forest
reduction}\label{constraint-graphs-and-forest-reduction}}

For each source color \texttt{a\ in\ O}, define a graph

\[
\Gamma_a=\Gamma_a(S)
\]

on vertex set \texttt{R} by

\[
\{i,j\}\in E(\Gamma_a)
\quad\Longleftrightarrow\quad
[i,j;a]\in S.
\tag{8}
\]

Equality propagates along paths. Hence only the connected-component
partition of \texttt{Gamma\_a} matters.

\hypertarget{lemma-3.1-forest-reduction}{%
\subsubsection{Lemma 3.1 --- forest
reduction}\label{lemma-3.1-forest-reduction}}

Every constraint system can be replaced, without changing its satisfying
tuples and without increasing its size, by one in which every
\texttt{Gamma\_a} is a forest.

\textbf{Proof.} Replace every connected component of \texttt{Gamma\_a}
by an arbitrary spanning tree. The new constraints impose exactly the
same equality relation on the values \texttt{pi\_i(a)}. Repeating
independently for all \texttt{a} proves the claim. \texttt{square}

For a reduced system let

\[
m_a=|E(\Gamma_a)|,
\qquad
c_a=\kappa(\Gamma_a)=r-m_a.
\tag{9}
\]

Then

\[
|S|=\sum_a m_a.
\tag{10}
\]

\begin{center}\rule{0.5\linewidth}{0.5pt}\end{center}

\hypertarget{synchronization-is-a-unique-coloring-problem}{%
\subsection{4. Synchronization is a unique-coloring
problem}\label{synchronization-is-a-unique-coloring-problem}}

Let \texttt{B\_a} be the set of connected components of
\texttt{Gamma\_a}. Construct a graph \texttt{H(S)} as follows.

Its vertices are

\[
(a,B),\qquad a\in O,\ B\in B_a.
\tag{11}
\]

For every phase index \texttt{i}, let \texttt{B\_a(i)} be the component
of \texttt{Gamma\_a} containing \texttt{i}. Join all \texttt{q} vertices

\[
(a,B_a(i)),\qquad a\in O,
\tag{12}
\]

pairwise. Thus every phase contributes a canonical \texttt{K\_q}
transversal. There is a canonical proper coloring

\[
\chi_0(a,B)=a.
\tag{13}
\]

\hypertarget{theorem-4.1-synchronization-unique-coloring-equivalence}{%
\subsubsection{Theorem 4.1 --- synchronization / unique-coloring
equivalence}\label{theorem-4.1-synchronization-unique-coloring-equivalence}}

A constraint family \texttt{S} is synchronizing if and only if the
canonical coloring \texttt{chi\_0} of \texttt{H(S)} is the unique proper
\texttt{q}-coloring up to permutation of the \texttt{q} color names.

\textbf{Proof.} Suppose first that \texttt{(pi\_1,...,pi\_r)} satisfies
\texttt{S}. If \texttt{i,j} lie in the same component \texttt{B} of
\texttt{Gamma\_a}, path propagation gives \texttt{pi\_i(a)=pi\_j(a)}.
Therefore

\[
F(a,B):=\pi_i(a)\qquad(i\in B)
\tag{14}
\]

is well defined. On the transversal (12) contributed by index
\texttt{i}, the \texttt{F}-values are

\[
\pi_i(0),\ldots,\pi_i(q-1),
\]

which are pairwise distinct. Hence \texttt{F} is a proper
\texttt{q}-coloring of \texttt{H(S)}.

Conversely, let \texttt{F} be a proper \texttt{q}-coloring. Every
canonical transversal is a \texttt{K\_q}, so its vertices receive all
\texttt{q} colors exactly once. Define

\[
\pi_i(a)=F(a,B_a(i)).
\tag{15}
\]

Then each \texttt{pi\_i} is a permutation. If
\texttt{{[}i,j;a{]}\ in\ S}, vertices \texttt{i,j} lie in the same
\texttt{Gamma\_a} component and (15) gives \texttt{pi\_i(a)=pi\_j(a)}.
Thus proper colorings are in bijection with satisfying phase tuples.

The diagonal tuples are exactly the colorings obtained from
\texttt{chi\_0} by one global output-color permutation. Hence
synchronization is equivalent to uniqueness of \texttt{chi\_0} up to
global relabeling. \texttt{square}

This theorem places the extremal problem next to classical unique
colorability, while preserving a severe additional restriction: the
quotient must arise from \texttt{r} canonical \texttt{K\_q} transversals
and identifications occur only within one source-color class.

\begin{center}\rule{0.5\linewidth}{0.5pt}\end{center}

\hypertarget{pair-union-connectivity}{%
\subsection{5. Pair-union connectivity}\label{pair-union-connectivity}}

\hypertarget{theorem-5.1}{%
\subsubsection{Theorem 5.1}\label{theorem-5.1}}

If \texttt{S} is synchronizing, then for every two distinct source
colors \texttt{a,b},

\[
\boxed{\Gamma_a\cup\Gamma_b\text{ is connected}.}
\tag{16}
\]

\textbf{Proof.} Suppose \texttt{Gamma\_a\ union\ Gamma\_b} is
disconnected. Let \texttt{X} be a nonempty proper union of its connected
components. No constraint of source color \texttt{a} or \texttt{b}
crosses the cut \texttt{(X,R\textbackslash{}X)}. Put
\texttt{sigma=(a\ b)} and define

\[
\pi_i=\sigma\quad(i\in X),
\qquad
\pi_i=id\quad(i\notin X).
\tag{17}
\]

Every internal constraint compares equal permutations. Every crossing
constraint has a source color different from \texttt{a,b}, hence fixed
by \texttt{sigma}. Thus all constraints are satisfied, but the tuple is
non-diagonal. Contradiction. \texttt{square}

For a reduced system, (16) implies

\[
m_a+m_b\ge r-1.
\tag{18}
\]

Summing (18) over all unordered pairs gives

\[
(q-1)\sum_a m_a\ge\binom q2(r-1).
\]

\hypertarget{corollary-5.2-half-density-bound}{%
\subsubsection{Corollary 5.2 --- half-density
bound}\label{corollary-5.2-half-density-bound}}

\[
\boxed{
L_q(r)\ge\left\lceil\frac{q(r-1)}2\right\rceil.
}
\tag{19}
\]

The connectivity statement is the LQR specialization of a classical
property of uniquely colorable graphs {[}3{]}. Its role here is as the
first lower-bound engine in the restricted transversal quotient.

\begin{center}\rule{0.5\linewidth}{0.5pt}\end{center}

\hypertarget{exact-low-dimensional-families}{%
\subsection{6. Exact low-dimensional
families}\label{exact-low-dimensional-families}}

\hypertarget{binary-alphabet}{%
\subsubsection{6.1 Binary alphabet}\label{binary-alphabet}}

For \texttt{q=2}, agreement on one source color identifies two
permutations of a two-element set. Therefore a connected graph on the
\texttt{r} phase indices is necessary and sufficient.

\hypertarget{theorem-6.1}{%
\subsubsection{Theorem 6.1}\label{theorem-6.1}}

\[
\boxed{L_2(r)=r-1.}
\tag{20}
\]

\hypertarget{two-phases}{%
\subsubsection{6.2 Two phases}\label{two-phases}}

Two permutations of a \texttt{q}-element set that agree on \texttt{q-1}
source points agree on the final point. Fewer than \texttt{q-1}
constraints leave at least two unconstrained source points and allow a
transposition there.

\hypertarget{theorem-6.2}{%
\subsubsection{Theorem 6.2}\label{theorem-6.2}}

\[
\boxed{L_q(2)=q-1.}
\tag{21}
\]

\hypertarget{three-colors}{%
\subsubsection{6.3 Three colors}\label{three-colors}}

For \texttt{q=3}, (18) gives

\[
m_0+m_1\ge r-1,
\quad
m_0+m_2\ge r-1,
\quad
m_1+m_2\ge r-1,
\]

and therefore

\[
L_3(r)\ge\left\lceil\frac{3(r-1)}2\right\rceil.
\tag{22}
\]

The matching construction is a three-constraint gadget. Suppose
\texttt{rho} is already synchronized and \texttt{sigma,tau\ in\ S\_3}
are new. Impose

\[
[\rho,\sigma;0],
\qquad
[\rho,\tau;1],
\qquad
[\sigma,\tau;2].
\tag{23}
\]

After global normalization take \texttt{rho=id}. Then \texttt{sigma}
fixes \texttt{0}, \texttt{tau} fixes \texttt{1}, and
\texttt{sigma(2)=tau(2)}. If \texttt{sigma} were the nontrivial
permutation fixing \texttt{0}, then \texttt{sigma(2)=1}, impossible for
\texttt{tau(2)} because \texttt{tau(1)=1}. Hence \texttt{sigma=id}; then
\texttt{tau(2)=2}, and \texttt{tau=id}.

Attaching two phases at a time gives the lower bound exactly.

\hypertarget{theorem-6.3}{%
\subsubsection{Theorem 6.3}\label{theorem-6.3}}

For every \texttt{r\textgreater{}=1},

\[
\boxed{
L_3(r)=\left\lceil\frac{3(r-1)}2\right\rceil.
}
\tag{24}
\]

\hypertarget{three-phases}{%
\subsubsection{6.4 Three phases}\label{three-phases}}

For \texttt{r=3}, every reduced \texttt{Gamma\_a} has
\texttt{m\_a\ in\ \{0,1,2\}}. If one color has \texttt{m\_a=0}, every
other color must be connected, giving cost at least \texttt{2(q-1)}.
Otherwise each color has at least one edge. Two one-edge colors must use
different edges of the phase triangle for their union to be connected.
There are only three edges, so at most three colors have
\texttt{m\_a=1}; all remaining colors have \texttt{m\_a=2}. Hence

\[
L_q(3)\ge3+2(q-3)=2q-3.
\tag{25}
\]

For the matching construction use the three-color triangle gadget

\[
[0,1;0],\quad[0,2;1],\quad[1,2;2]
\tag{26}
\]

and, for every extra source color \texttt{a\textgreater{}=3}, impose
\texttt{{[}0,1;a{]}} and \texttt{{[}0,2;a{]}}. The extra colors are
fixed across all phases, reducing the remaining action to \texttt{S\_3},
where (26) synchronizes.

\hypertarget{theorem-6.4}{%
\subsubsection{Theorem 6.4}\label{theorem-6.4}}

For every \texttt{q\textgreater{}=3},

\[
\boxed{L_q(3)=2q-3.}
\tag{27}
\]

\begin{center}\rule{0.5\linewidth}{0.5pt}\end{center}

\hypertarget{the-exact-four-phase-column}{%
\subsection{7. The exact four-phase
column}\label{the-exact-four-phase-column}}

For \texttt{r=4}, every source color determines a partition
\texttt{P\_a} of four phase indices. Under forest reduction its rank

\[
m_a=4-|P_a|
\]

lies in \texttt{\{0,1,2,3\}}.

The compatibility condition from Theorem 5.1 becomes

\[
P_a\vee P_b=\mathbf 1
\qquad(a\ne b),
\tag{28}
\]

where \texttt{mathbf\ 1} is the one-block partition.

The four-point partition lattice gives the complete low-rank
classification:

\begin{itemize}
\tightlist
\item
  the discrete rank-zero partition is compatible only with the one-block
  rank-three partition;
\item
  two rank-one partitions are never compatible;
\item
  a fixed rank-one partition is compatible with exactly four rank-two
  partitions;
\item
  there are exactly seven rank-two partitions, and any two distinct ones
  are compatible;
\item
  a rank-three partition is compatible with all partitions.
\end{itemize}

Let \texttt{n\_k} be the number of colors of rank \texttt{k}. If no
rank-zero color occurs,

\[
|S|=n_1+2n_2+3n_3=3q-(2n_1+n_2).
\tag{29}
\]

The classification yields the maximal defect

\[
2n_1+n_2=\begin{cases}
q+1,&3\le q\le5,\\
6,&q=6,\\
7,&q\ge7.
\end{cases}
\tag{30}
\]

and therefore the corresponding lower bounds.

For \texttt{q=3}, the five-constraint construction is exactly an
instance of the already proved \texttt{S\_3} recursive gadget in Theorem
6.3, after phase/color relabeling. For \texttt{q=4,5}, one may use one
rank-one partition plus respectively three or four of its compatible
rank-two partners; direct normalization and injectivity force all phases
to the identity. For \texttt{q=6}, extend an optimal \texttt{q=5} gadget
by making the sixth source color connected. For \texttt{q=7}, use all
seven bipartitions of four phase indices; this is the \texttt{r=4}
instance of the binary-cut construction proved in Section 8. Extra
colors for \texttt{q\textgreater{}7} are made connected.

\hypertarget{theorem-7.1-exact-four-phase-cost}{%
\subsubsection{Theorem 7.1 --- exact four-phase
cost}\label{theorem-7.1-exact-four-phase-cost}}

\[
\boxed{
L_q(4)=
\begin{cases}
3,&q=2,\\
2q-1,&3\le q\le5,\\
12,&q=6,\\
3q-7,&q\ge7.
\end{cases}
}
\tag{31}
\]

The isolated value at \texttt{q=6} is therefore structural rather than
numerical noise: the rank-one architecture can coexist with only four
rank-two colors, whereas the all-bipartition architecture has capacity
seven.

\begin{center}\rule{0.5\linewidth}{0.5pt}\end{center}

\hypertarget{universal-binary-cut-synchronization}{%
\subsection{8. Universal binary-cut
synchronization}\label{universal-binary-cut-synchronization}}

Fix \texttt{r\textgreater{}=2} and put

\[
n=r-1,
\qquad
V=\mathbb F_2^n\setminus\{0\}.
\tag{32}
\]

Use one active source color for each vector

\[
v=(v_1,\ldots,v_n)\in V.
\]

Thus

\[
q_0=|V|=2^{r-1}-1.
\tag{33}
\]

For color \texttt{v}, partition the phases into

\[
B_v^0=\{0\}\cup\{i:v_i=0\},
\qquad
B_v^1=\{i:v_i=1\}.
\tag{34}
\]

Connect each block internally by a spanning tree. Since there are two
nonempty blocks, color \texttt{v} costs \texttt{r-2} constraints.

\hypertarget{theorem-8.1-binary-cut-gadget}{%
\subsubsection{Theorem 8.1 --- binary-cut
gadget}\label{theorem-8.1-binary-cut-gadget}}

The family (34) synchronizes the \texttt{r} phases in
\texttt{S\_\{2\^{}\{r-1\}-1\}}.

\textbf{Proof.} Normalize \texttt{pi\_0=id}. Fix
\texttt{i\textgreater{}=1} and define

\[
H_i=\{v:v_i=0\},
\qquad
A_i=\{v:v_i=1\}.
\tag{35}
\]

If \texttt{v\ in\ H\_i}, phases \texttt{0} and \texttt{i} lie in the
same connected block \texttt{B\_v\^{}0}, so

\[
\pi_i(v)=v.
\tag{36}
\]

Thus \texttt{pi\_i} fixes \texttt{H\_i} pointwise and preserves
\texttt{A\_i} setwise.

For \texttt{j!=i} and \texttt{v\ in\ A\_i\ cap\ A\_j}, phases
\texttt{i,j} lie in the same block \texttt{B\_v\^{}1}, hence

\[
\pi_i(v)=\pi_j(v).
\tag{37}
\]

The left side belongs to \texttt{A\_i} and the right side to
\texttt{A\_j}; therefore the common value lies in
\texttt{A\_i\ cap\ A\_j}. Hence \texttt{pi\_i} preserves every subset
\texttt{A\_i\ cap\ A\_j} setwise.

On \texttt{A\_i}, the membership bits in the sets
\texttt{A\_i\ cap\ A\_j} for \texttt{j!=i} are precisely the remaining
coordinates \texttt{(v\_j)\_\{j!=i\}}. They distinguish all vectors of
\texttt{A\_i}. Thus \texttt{pi\_i} fixes every vector of \texttt{A\_i},
and together with (36) this gives \texttt{pi\_i=id}. This holds for all
\texttt{i}, so all phases are diagonal before normalization.
\texttt{square}

The gadget costs

\[
(r-2)(2^{r-1}-1).
\tag{38}
\]

For every \texttt{q\textgreater{}q\_0}, make each extra source color
connected on the \texttt{r} phases at cost \texttt{r-1}. Therefore

\[
L_q(r)\le(r-1)q-(2^{r-1}-1)
\quad(q\ge q_0).
\tag{39}
\]

\begin{center}\rule{0.5\linewidth}{0.5pt}\end{center}

\hypertarget{cut-spaces-and-the-packing-bound}{%
\subsection{9. Cut spaces and the packing
bound}\label{cut-spaces-and-the-packing-bound}}

The matching lower bound comes from a canonical encoding of every
source-color component partition.

Let \texttt{P} be a partition of the phase set

\[
R=\{0,1,\ldots,r-1\}.
\]

Normalize binary cuts by fixing the bit at phase \texttt{0} to zero and
identify them with

\[
\mathbb F_2^{r-1}.
\]

Define

\[
W(P)=\{x\in\mathbb F_2^{r-1}:x\text{ is constant on every block of }P\},
\tag{40}
\]

with the understood extension \texttt{x\_0=0}.

\hypertarget{lemma-9.1}{%
\subsubsection{Lemma 9.1}\label{lemma-9.1}}

If \texttt{P} has \texttt{c} blocks, then

\[
\dim W(P)=c-1.
\tag{41}
\]

\textbf{Proof.} The block containing phase \texttt{0} is forced to bit
zero. Every other block chooses its bit independently. Hence
\texttt{\textbar{}W(P)\textbar{}=2\^{}\{c-1\}}. \texttt{square}

For the component partition \texttt{P\_a} of a forest-reduced
\texttt{Gamma\_a}, define the saved cost

\[
d_a=(r-1)-m_a.
\tag{42}
\]

Since \texttt{m\_a=r-c\_a},

\[
d_a=c_a-1=\dim W(P_a).
\tag{43}
\]

\hypertarget{lemma-9.2-joinintersection-identity}{%
\subsubsection{Lemma 9.2 --- join/intersection
identity}\label{lemma-9.2-joinintersection-identity}}

For any partitions \texttt{P,Q},

\[
\boxed{W(P)\cap W(Q)=W(P\vee Q).}
\tag{44}
\]

\textbf{Proof.} A normalized binary cut belongs to both spaces exactly
when it is constant on every block of \texttt{P} and every block of
\texttt{Q}. This is equivalent to being constant on every equivalence
class generated jointly by those blocks, namely the blocks of
\texttt{P\ vee\ Q}. \texttt{square}

By Theorem 5.1,

\[
P_a\vee P_b=\mathbf 1
\quad(a\ne b),
\]

so

\[
W(P_a)\cap W(P_b)=\{0\}.
\tag{45}
\]

Thus the nonzero vectors of the cut spaces are pairwise disjoint. Since
\texttt{F\_2\^{}\{r-1\}} has \texttt{2\^{}\{r-1\}-1} nonzero vectors:

\hypertarget{theorem-9.3-defect-packing-inequality}{%
\subsubsection{Theorem 9.3 --- defect packing
inequality}\label{theorem-9.3-defect-packing-inequality}}

Every synchronizing system satisfies

\[
\boxed{
\sum_a(2^{d_a}-1)\le2^{r-1}-1.
}
\tag{46}
\]

The counting step in (46) is standard finite-geometry packing once the
cut-space reduction has been made; it is not claimed here as a new
theorem about partial spreads or vector-space partitions {[}5,6{]}.

Since

\[
d\le2^d-1
\qquad(d\ge0),
\tag{47}
\]

we get

\[
\sum_a d_a\le2^{r-1}-1.
\tag{48}
\]

But by (42),

\[
|S|=(r-1)q-\sum_a d_a.
\tag{49}
\]

Therefore:

\hypertarget{corollary-9.4-universal-defect-lower-bound}{%
\subsubsection{Corollary 9.4 --- universal defect lower
bound}\label{corollary-9.4-universal-defect-lower-bound}}

\[
\boxed{
L_q(r)\ge(r-1)q-(2^{r-1}-1).
}
\tag{50}
\]

\begin{center}\rule{0.5\linewidth}{0.5pt}\end{center}

\hypertarget{exact-large-alphabet-stabilization}{%
\subsection{10. Exact large-alphabet
stabilization}\label{exact-large-alphabet-stabilization}}

We now combine Sections 8 and 9.

\hypertarget{theorem-10.1-exact-stabilization-theorem}{%
\subsubsection{Theorem 10.1 --- exact stabilization
theorem}\label{theorem-10.1-exact-stabilization-theorem}}

For every \texttt{r\textgreater{}=2} and every

\[
q\ge2^{r-1}-1,
\]

\[
\boxed{
L_q(r)=(r-1)q-(2^{r-1}-1).
}
\tag{51}
\]

Moreover the threshold is exact: if

\[
q<2^{r-1}-1,
\]

then

\[
L_q(r)>(r-1)q-(2^{r-1}-1).
\tag{52}
\]

\textbf{Proof.} The lower bound is Corollary 9.4. The binary-cut
construction gives equality when
\texttt{q\textgreater{}=2\^{}\{r-1\}-1}.

Suppose equality held in (50). Then

\[
\sum_a d_a=2^{r-1}-1.
\tag{53}
\]

From (46) and (47),

\[
\sum_a d_a
\le
\sum_a(2^{d_a}-1)
\le
2^{r-1}-1.
\tag{54}
\]

Equality in (53) forces equality termwise in (47), which for nonnegative
integers occurs only for \texttt{d\_a\ in\ \{0,1\}}. Hence every
positive-defect color contributes exactly one nonzero cut vector. To
obtain total defect \texttt{2\^{}\{r-1\}-1}, at least that many
positive-defect colors are necessary. Therefore

\[
q\ge2^{r-1}-1.
\]

So equality is impossible below that threshold. \texttt{square}

\hypertarget{corollary-10.2}{%
\subsubsection{Corollary 10.2}\label{corollary-10.2}}

The stabilization threshold is

\[
\boxed{q_0(r)=2^{r-1}-1.}
\tag{55}
\]

For example,

\[
L_q(5)=4q-15\qquad(q\ge15),
\tag{56}
\]

and

\[
L_q(6)=5q-31\qquad(q\ge31).
\tag{57}
\]

Thus for every fixed \texttt{r}, only finitely many alphabet sizes
remain outside the exact linear tail.

\begin{center}\rule{0.5\linewidth}{0.5pt}\end{center}

\hypertarget{the-finite-pre-stabilization-problem}{%
\subsection{11. The finite pre-stabilization
problem}\label{the-finite-pre-stabilization-problem}}

The proof of Theorem 10.1 gives a stronger interface for the unresolved
sector. Define

\[
\mathcal D_r(q)
=
\max\left\{
\sum_{i=1}^q\dim W(P_i):
W(P_i)\cap W(P_j)=\{0\}\ (i\ne j)
\right\},
\tag{58}
\]

where zero-dimensional spaces may pad the family.

Then necessarily

\[
L_q(r)\ge(r-1)q-\mathcal D_r(q).
\tag{59}
\]

This is only a lower-bound interface: pairwise trivial intersection is
necessary for synchronization but need not by itself imply unique
colorability of the transversal quotient.

The genuinely unresolved regime is

\[
\boxed{
4\le q<2^{r-1}-1,
\qquad r\ge5.
}
\tag{60}
\]

For \texttt{r=5}, only \texttt{4\textless{}=q\textless{}=14} remain.
Exact weighted packing computations provide a useful hostile check, but
no unproved packing optimum is promoted to an LQR theorem here.

\begin{center}\rule{0.5\linewidth}{0.5pt}\end{center}

\hypertarget{finite-verification}{%
\subsection{12. Finite verification}\label{finite-verification}}

Two independent verification scripts accompany the theorem development
in the repository.

\texttt{verify\_lqr.py} checks:

\begin{itemize}
\tightlist
\item
  the exact \texttt{q=3} constructions for small \texttt{r} by
  normalized exhaustive enumeration;
\item
  the exact \texttt{r=3} constructions for small \texttt{q};
\item
  the first exact \texttt{r=4} values using the partition quotient.
\end{itemize}

\texttt{verify\_lqr\_cutspace.py} independently checks:

\begin{itemize}
\tightlist
\item
  Bell-number enumeration of phase partitions through \texttt{r=6};
\item
  \texttt{\textbar{}W(P)\textbar{}=2\^{}\{\textbar{}P\textbar{}-1\}};
\item
  \texttt{W(P)\ cap\ W(Q)=W(P\ vee\ Q)} for all partition pairs through
  \texttt{r=6};
\item
  the weighted partition-packing capacity for \texttt{r=5} through the
  stabilization threshold \texttt{q=15}.
\end{itemize}

The proofs of the infinite formulas do not depend on these computations.

\begin{center}\rule{0.5\linewidth}{0.5pt}\end{center}

\hypertarget{literature-position-and-novelty-boundary}{%
\subsection{13. Literature position and novelty
boundary}\label{literature-position-and-novelty-boundary}}

The graph-theoretic and finite-geometric surroundings of the result are
classical.

Harary, Hedetniemi and Robinson {[}3{]} established foundational
properties of uniquely colorable graphs, including connectivity of the
graph induced by every pair of color classes. Bollobas {[}4{]} developed
further structural results on unique colorability. Our Theorem 5.1 is
therefore not claimed as a new general unique-colorability theorem; its
role is the specialized form induced by the LQR transversal quotient.

Pairwise trivially intersecting subspaces, partial spreads and
vector-space partitions are also classical {[}5,6{]}. Once (45) has been
obtained, counting their nonzero vectors is standard. The novelty claim
is therefore deliberately restricted to the following chain inside the
FCOA phase-synchronization problem:

\[
\text{point-image constraints}
\longrightarrow
\text{transversal unique-coloring quotient}
\longrightarrow
\text{component partitions}
\longrightarrow
\text{canonical binary cut spaces}
\longrightarrow
\text{sharp LQR stabilization}.
\tag{61}
\]

Permutation synchronization also has an established algorithmic
literature, typically concerned with recovering globally consistent
matchings from noisy pairwise permutation estimates {[}7{]}. The present
parameter differs: it asks for the minimum number of exact same-source
point-image equalities sufficient to force an arbitrary tuple in
\texttt{S\_q\^{}r} to be diagonal.

A dedicated search performed for this project did not locate the exact
extremal parameter \texttt{L\_q(r)} in this form. This is a conservative
literature statement, not an absolute priority claim.

\begin{center}\rule{0.5\linewidth}{0.5pt}\end{center}

\hypertarget{limitations-and-open-problems}{%
\subsection{14. Limitations and open
problems}\label{limitations-and-open-problems}}

The main limitations are structural and explicit.

\begin{enumerate}
\def\labelenumi{\arabic{enumi}.}
\tightlist
\item
  \texttt{L\_q(r)} measures abstract phase-link constraints, not real
  FCOA operation-cell additions.
\item
  The pre-stabilization sector (60) is not solved in general.
\item
  Pairwise cut-space packing is necessary but may fail to capture all
  unique-coloring obstructions.
\item
  The present paper does not define a multicolor analogue of the binary
  actual-cell repair parameter.
\item
  The local existence of permutation-valued phases is assumed; sparse
  ternary reducts with \texttt{q\textgreater{}=3} can fail before the
  synchronization problem is even defined, as shown in the companion
  sparse multicolor transport manuscript {[}2{]}.
\end{enumerate}

Natural next problems are to determine \texttt{L\_4(r)} for general
\texttt{r}, solve \texttt{L\_q(5)} for
\texttt{4\textless{}=q\textless{}=14}, and characterize when an optimal
partition-subspace packing is realized by a synchronizing transversal
quotient.

\begin{center}\rule{0.5\linewidth}{0.5pt}\end{center}

\hypertarget{conclusion}{%
\subsection{15. Conclusion}\label{conclusion}}

The point-image synchronization cost has a sharper structure than the
naive spanning-tree estimate suggests. The correct lower-bound objects
are not merely graphs on the phase indices. They are component
partitions, and after normalization those partitions produce binary cut
spaces whose nonzero vectors must pack disjointly.

This leads to three layers of exact structure:

\[
L_3(r)=\left\lceil\frac{3(r-1)}2\right\rceil,
\]

an exact four-phase transition,

\[
L_q(4)=3,\ 2q-1,\ 12,\ 3q-7
\]

in its respective ranges, and the general stabilization theorem

\[
\boxed{
L_q(r)=(r-1)q-(2^{r-1}-1)
\quad\text{iff the linear lower bound is attainable, with attainment beginning exactly at }q=2^{r-1}-1.
}
\]

For each fixed number of phases, the unsolved problem is therefore
finite. In FCOA terms, the result provides an exact capacity theorem for
synchronizing already-existing anonymous full-support phases, while
keeping distinct the harder problem of realizing such synchronization by
genuine operation-cell extensions.

\begin{center}\rule{0.5\linewidth}{0.5pt}\end{center}

\hypertarget{references}{%
\subsection{References}\label{references}}

{[}1{]} A. Malachevsky and Commander Sol, \textbf{Fixed-Carrier Oriented
Algebra (FCOA): Definition, Typed Partial Operations, Carrier Erasure,
and the Canonical M0 Baseline}, 2026. DOI:
\texttt{10.5281/zenodo.22164246}.
https://doi.org/10.5281/zenodo.22164246.

{[}2{]} A. Malachevsky and Commander Sol, \textbf{Reflections on Sparse
Multicolor Transport with Commander Sol: Proper-Coloring Obstructions,
Local Phase Groupoids, and Exact Gluing in FCOA}, companion manuscript,
2026, repository package
\texttt{papers/FCOA-SPARSE-MULTICOLOR-TRANSPORT/}.

{[}3{]} F. Harary, S. T. Hedetniemi, R. W. Robinson, \textbf{Uniquely
colorable graphs}, \emph{Journal of Combinatorial Theory} 6(3) (1969),
264--270. DOI: \texttt{10.1016/S0021-9800(69)80086-4}.

{[}4{]} B. Bollobas, \textbf{Uniquely colorable graphs}, \emph{Journal
of Combinatorial Theory, Series B} 25 (1978), 54--61. DOI:
\texttt{10.1016/S0095-8956(78)80010-0}.

{[}5{]} O. Heden, \textbf{A survey of the different types of vector
space partitions}, \emph{Discrete Mathematics, Algorithms and
Applications} 4(1) (2012), 1250001. DOI:
\texttt{10.1142/S1793830912500012}.

{[}6{]} T. Honold, M. Kiermaier, S. Kurz, \textbf{Partial spreads and
vector space partitions}, in \emph{Network Coding and Subspace Designs},
Springer, 2018, 131--170. DOI: \texttt{10.1007/978-3-319-70293-3\_7}.

{[}7{]} D. Pachauri, T. Kondor, V. Singh, \textbf{Solving the multi-way
matching problem by permutation synchronization}, \emph{Advances in
Neural Information Processing Systems 26}, 2013.

\begin{center}\rule{0.5\linewidth}{0.5pt}\end{center}

\hypertarget{reproducibility-note}{%
\subsection{Reproducibility note}\label{reproducibility-note}}

The theorem-source files and independent finite verifiers are maintained
under

\texttt{delegated/FCOA\_RIGIDITY\_COST/QGE3/}

in the repository

\texttt{https://github.com/AIDevelopersMonster/Riemann-Hypothesis-Commander-Sol}

on the research branch \texttt{director/fcoa-rigidity-cost} at
manuscript assembly time.
